\documentclass{article}
\usepackage{/globals/de}
\usepackage{/globals/code}
\usepackage{/globals/anki}
\ankiprefix{sql}
\ankitopic{SQL}
\ankishow
\begin{document}
\section{SQL} 
Mit SQL, der Structured Query Language, werden Datenbanken verwaltet.

Befehle enden mit einem Semikolon, Keywörter werden groß geschrieben. Namen mit Leerzeichen werden in eckigen Klammern geschrieben.
\ankinote{leerzeichen}{Namen mit Leerzeichen}{Eckige Klammern}

\subsection{\ilsourcecode[sql]|CREATE|}
Mit \ilsourcecode[sql]|CREATE| können Datenbanken, Tabellen, Indexe, \etc erstellt werden.
\begin{sourcecode}{sql}
CREATE DATABASE Datenbankname;
\end{sourcecode}
\begin{sourcecode}{sql}
CREATE TABLE Tabellenname;
\end{sourcecode}
Bei der Erstellung von Tabellen können auch direkt die Spalten angegeben werden.
\begin{sourcecode}{sql}
CREATE DATABASE Tabellenname (
    Spaltenname1 datenwert1,
    Spaltenname2 datenwert2
);
\end{sourcecode}
\ankinote{erstellen}{Create}{CREATE DATABASE Tabellenname (Spaltenname1 datenwert1, Spaltenname2 datenwert2);}

\subsection{\ilsourcecode[sql]|INSERT|}
Mit \ilsourcecode[sql]|INSERT| können neue Datensätze hinzugefügt werden. Dabei werden die zu füllenden Spalten genannt, dann, in der gleichen Reihenfolge, dessen Werte.
\begin{sourcecode}{sql}
INSERT INTO Tabellenname (Spaltenname1, Spaltenname2) VALUES (wert1, wert2);
\end{sourcecode}
Wenn alle Spalten Werte haben müssen die Spaltennamen nicht wiederholt werden.
\begin{sourcecode}{sql}
INSERT INTO Tabellenname VALUES (wert1, wert2);
\end{sourcecode}
\ankinote{insert-short}{Insert kurz}{INSERT INTO Tabellenname VALUES (wert1, wert2);}
\ankinote{insert-full}{Insert voll}{INSERT INTO Tabellenname (Spaltenname1, Spaltenname2) VALUES (wert1, wert2);}

\subsection{\ilsourcecode[sql]|SELECT|}
Mit \ilsourcecode[sql]|SELECT| können Daten ausgelesen werden. Ein \ilsourcecode[sql]|*| steht für alle Spalten.
\begin{sourcecode}{sql}
SELECT * FROM Tabellenname;
\end{sourcecode}
Ansonsten können auch einzelne Spalten namentlich angegeben werden.
\begin{sourcecode}{sql}
SELECT Spaltenname1, Spaltenname2 FROM Tabellenname;
\end{sourcecode}
Es kann nach Datensätzen die bestimmte Konditionen erfüllen gesucht werden.
\begin{sourcecode}{sql}
SELECT * FROM Tabellenname WHERE Spaltenname=wert;
\end{sourcecode}
Datensätze können anhand von einem Wert sotiert werden, aufsteigend (\ilsourcecode[sql]|ASC|) oder absinkend (\ilsourcecode[sql]|DESC|).
\begin{sourcecode}{sql}
SELECT * FROM Tabellenname ORDER BY Spaltenname ASC;
\end{sourcecode}
\ankinote{order}{Oder}{aufsteigend ASC, absinkend DESC}
Es können auch Funktionen über ganze Datensätze laufen.
\begin{sourcecode}{sql}
SELECT COUNT(*) FROM Tabellenname;
\end{sourcecode}
Gleiches gilt für \ilsourcecode[sql]{AVG(Spaltenname)}, \ilsourcecode[sql]{SUM(Spaltenname)}, \ilsourcecode[sql]{MIN(Spaltenname)} und \ilsourcecode[sql]{MAX(Spaltenname)}.
\ankinote{funktionen}{Funktionen}{COUNT, AVG, SUM, MIN, MAX}
\end{document}